% Capitolo da inserire fra "Sicurezza e distanze nella navigazione sociale" e "Correzione
% neurale della previsione di moto". I capitoli successivi slittano di uno.
%
% Riscontri nel codice (main.py):
%   2155  algoritmo_a_star(griglia_robot, ...)      il robot cerca sulla griglia fine
%   1118  _pool_worker_griglia da _celle_muro_di(griglia)   la folla cerca su quella grossa
%   1995  celle bloccate del lidar sulla griglia fine       (varchi fra persone vicine)
%   2109  rapporto_collasso = sottocelle_per_lato // fattore_robot
%    309  costo_occupata = COSTO_CELLA_OCCUPATA * dim / DIM_NODO
%   2110  costo_massimo_probabilita ... * dim_robot / DIM_NODO
%   2132  costo_isteresi ... * dim_robot / DIM_NODO
%   2032  costo_alone ... * dim_robot / DIM_NODO

\section{Ottimizzazione della risoluzione delle griglie}
\label{cap:ottimizzazione-griglie}

I due capitoli precedenti hanno presentato la pianificazione e la rappresentazione del rischio
come se entrambe insistessero sulla medesima discretizzazione dello spazio. Non è così, ed è una
scelta deliberata: il sistema impiega tre risoluzioni distinte, ciascuna dimensionata sul
requisito della componente che la utilizza.

La ragione è che la risoluzione della griglia governa contemporaneamente due grandezze che si
muovono in versi opposti. Da un lato la fedeltà con cui si rappresenta la geometria — la forma
di una regione di rischio, l'esistenza di un varco fra due persone vicine — che migliora al
diminuire del passo. Dall'altro il costo della ricerca del percorso, che cresce
approssimativamente con il cubo dell'inverso del passo: dimezzare il lato della cella
quadruplica i nodi da espandere e raddoppia i campioni richiesti da ciascun controllo di
visibilità, operazione che assorbe la quasi totalità del tempo di pianificazione. Una griglia
unica costringerebbe a scegliere un compromesso valido per tutti gli usi, e quindi ottimo per
nessuno.

\begin{table}[h]
\centering
\begin{tabular}{lcl}
\hline
griglia & lato & impiego \\
\hline
ambiente & $\ell_0$ & muri, occupazione, pianificazione della folla \\
ricerca del robot & $\ell_0/2$ & espansione dei nodi e controllo di visibilità del robot \\
regione di rischio & $\ell_0/2k$ & propagazione delle macchie di probabilità \\
\hline
\end{tabular}
\caption{Le tre risoluzioni impiegate. Ciascuna è un multiplo intero della successiva, così che
il passaggio dall'una all'altra sia una divisione priva di resto.}
\label{tab:tre-risoluzioni}
\end{table}

L'asimmetria fra robot e folla merita una parola. Le persone continuano a pianificare sulla
griglia dell'ambiente, con il medesimo algoritmo any-angle impiegato dal robot ma senza alcun
raffinamento. La scelta non è una semplificazione di comodo: la folla costituisce l'ambiente
rispetto al quale il robot viene misurato, e raffinarne la pianificazione ne aumenterebbe il
costo — che è già dominante, essendo le persone alcune centinaia contro un robot solo — senza
incidere sulla grandezza in esame. Il raffinamento è speso dove se ne osserva l'effetto.

Perché le tre risoluzioni possano coesistere deve però valere una condizione, che attraversa
l'intero capitolo: \textbf{i costi non devono dipendere dalla risoluzione}. Se la penalità
associata a una cella restasse la stessa al dimezzarsi del suo lato, attraversare una data
distanza costerebbe il doppio, e raffinare la griglia equivarrebbe ad aumentare di nascosto la
prudenza del robot. Ogni penalità è pertanto riscalata con il lato della cella cui è applicata,
senza eccezioni: la cella occupata, la regione di rischio, lo sconto di isteresi sul percorso
già intrapreso e l'alone attorno alle persone ferme. È questa invarianza a rendere la
risoluzione un parametro libero anziché una manopola implicita sul comportamento.

\subsection{Risoluzione della ricerca: la griglia del robot}
\label{sec:griglia-robot}

Il robot cerca il percorso su una griglia di lato dimezzato rispetto a quella dell'ambiente,
ottenuta suddividendo ciascuna cella in quattro e replicandone il tipo: i muri sono gli stessi,
soltanto descritti più finemente. Il guadagno non sta dunque nella geometria delle pareti, che è
invariata, bensì nella granularità di tutto ciò che sulla griglia viene appoggiato.

Il beneficio più rilevante riguarda l'occupazione. Le persone rilevate dal telemetro vengono
marcate come celle non percorribili, e su una griglia il cui passo è confrontabile con
l'ingombro di un pedone due persone vicine ma non a contatto occupano celle adiacenti,
sigillando un passaggio che nello spazio continuo resta aperto. Sulla griglia dimezzata le
medesime due persone lasciano fra sé celle libere, e il varco diventa percorribile. La
differenza non è estetica: in un ambiente affollato la maggior parte delle alternative di
percorso si gioca su passaggi di questa ampiezza, e sigillarli equivale a imporre al
pianificatore deviazioni che la geometria reale non richiede.

Sulla medesima griglia insistono inoltre l'alone di rispetto attorno alle persone ferme e lo
sconto di isteresi applicato al percorso già intrapreso. Entrambi sono campi definiti su
distanze dell'ordine del metro, e su celle di lato paragonabile risulterebbero quantizzati al
punto da perdere il proprio significato: un alone che valga una sola cella non descrive una
distanza di rispetto, descrive un ostacolo.

Il raffinamento non migliora, per contro, la geometria del percorso. Il pianificatore è
any-angle: produce tratti rettilinei a inclinazione arbitraria, non vincolati agli orientamenti
della griglia, e questa proprietà non dipende dal passo. Ciò che cambia con la risoluzione è
dove il percorso può passare — quali varchi esistono, quali celle costano — non con quale
angolazione lo attraversa.

\subsection{Risoluzione del rischio: la sotto-griglia delle macchie}
\label{sec:sottogriglia-rischio}

La regione di rischio pone un requisito ulteriore, e di natura diversa. Non è un insieme di
celle occupate ma un oggetto geometrico: si allunga lungo la direzione di marcia, si piega
attorno agli spigoli e si restringe in corrispondenza dei varchi. Ciascuna di queste proprietà è
rappresentabile soltanto se la regione è descritta da un numero sufficiente di celle; su una
griglia dal passo confrontabile con l'ingombro di una persona l'anisotropia degenera e la
deformazione attorno a un angolo si riduce a un gradino.

La macchia viene perciò propagata su una sotto-griglia di lato $\ell_0/2k$, con $k$ intero, e il
risultato aggregato sulla griglia di ricerca prima che il pianificatore lo consumi. Il valore di
$k$ è vincolato a rendere l'aggregazione una divisione intera, cosicché nessuna sotto-cella
ricada a cavallo di due celle di decisione.

Perché tale separazione sia legittima devono però valere tre proprietà, verificate nel seguito:
che il raffinamento non alteri la regione che descrive, che l'aggregazione non introduca a sua
volta una dipendenza dalla risoluzione, e che il costo aggiuntivo non ricada sulla componente
che già domina il bilancio di calcolo.

\subsubsection{Invarianza della forma}

Il raffinamento deve modificare la fedeltà con cui la regione è descritta, non la regione
stessa: diversamente $k$ diventerebbe una seconda manopola sull'estensione del rischio. Si è
pertanto imposto che il raggio di propagazione, espresso in sotto-celle, scali con $k$, e che il
decadimento per sotto-cella valga $\delta^{1/k}$. La seconda condizione discende dalla proprietà
moltiplicativa del peso: essendo questo il prodotto dei decadimenti dei passi che vi conducono,
elevare il fattore all'inverso di $k$ compensa esattamente il numero $k$ volte maggiore di passi
necessari a coprire la medesima distanza. Estensione fisica e decadimento per unità di lunghezza
restano dunque immutati: raffinando si ottiene la stessa macchia descritta meglio, non una
macchia diversa.

\subsubsection{Aggregazione per massimo}

A ciascuna cella di ricerca è attribuito il massimo fra i pesi delle sotto-celle che la
compongono. La somma renderebbe la penalità dipendente da $k$, crescendo il numero di addendi
con $k^2$, e violerebbe l'invarianza posta in apertura. La media diluirebbe un picco genuino nel
volume della cella che lo contiene: una sotto-cella a rischio elevato circondata da sotto-celle
libere produrrebbe una cella di costo modesto, e il pianificatore l'attraverserebbe. Il massimo
conserva invece la lettura conservativa — il punto peggiore che il robot incontrerebbe
attraversando quella cella — ed è indipendente dalla risoluzione.

\subsubsection{Località del costo}

A differenza del raffinamento della ricerca, che è uniforme su tutta la mappa, questo è locale
per costruzione: la propagazione si sviluppa attorno a ciascuna persona rilevata, entro un
raggio limitato e con interruzione al di sotto di una soglia di peso trascurabile. Il suo costo
è pertanto proporzionale al numero di persone percepite e non all'estensione dell'ambiente, e
non ricade sulla ricerca del percorso, che continua a operare sulla propria griglia e a
espandere il medesimo numero di nodi. Il raffinamento del rischio si paga soltanto dove un
rischio esiste; lo spazio libero, che costituisce la maggior parte della mappa, resta descritto
alla risoluzione grossolana che al pianificatore è sufficiente.

\subsection{Limiti del raffinamento e configurazione impiegata}
\label{sec:limiti-raffinamento}

Le due sezioni precedenti hanno esposto ciò che il raffinamento consente. Questa espone il
rovescio: perché non lo si estende a tutto lo spazio, quale porzione del beneficio
l'aggregazione trattiene, e in quale configurazione il sistema è stato effettivamente misurato.

La via apparentemente ovvia — portare tutto alla risoluzione più fine — è preclusa dal costo.
L'onere della ricerca cresce all'incirca con il cubo dell'inverso del passo, e il profilo di
esecuzione misurato su uno scenario denso colloca il 97\% del tempo di calcolo entro la ricerca
di percorso e il 96\% di quest'ultimo entro il solo controllo di visibilità, invocato
dell'ordine di venti milioni di volte in trecento fotogrammi. Il dimezzamento già adottato costa
quindi, di per sé, circa otto volte la pianificazione sulla griglia dell'ambiente: stima
ricavata dal prodotto fra il quadruplicarsi dei nodi e il raddoppiarsi dei campioni per
controllo, non da una misura diretta. Un ulteriore dimezzamento porterebbe l'onere a due ordini
di grandezza sopra la configurazione di partenza, a carico della componente che già domina il
bilancio, e sarebbe incompatibile tanto con le risorse di un sistema a bordo quanto con
l'esecuzione di campagne di validazione composte da centinaia di traversate appaiate.

Il raffinamento del rischio, dal canto suo, produce una rappresentazione più fedele ma non
abbassa la risoluzione alla quale il pianificatore decide: l'aggregazione riporta la macchia
sulla griglia di ricerca, ed è lì che la decisione avviene.
L'ottimizzazione migliora dunque la qualità dell'informazione trasmessa, non la finezza con cui
essa può essere utilizzata. La distinzione risulterà determinante nell'analisi condotta in
\S\ref{sec:dove-si-arresta}, dove si mostra che il guadagno di un'informazione più accurata si
trasferisce alla decisione soltanto nella misura in cui supera la risoluzione alla quale la
decisione è presa — soglia che il raffinamento della sola rappresentazione non sposta.

Restano acquisiti, entro tali limiti, due benefici indipendenti dall'accuratezza della
previsione e quindi non soggetti al vincolo appena descritto: l'apertura dei varchi fra corpi
vicini, che amplia l'insieme dei percorsi disponibili in ambiente affollato, e la fedeltà
geometrica delle regioni di rischio, che rende effettive l'anisotropia e la deformazione
attorno agli ostacoli descritte nel capitolo precedente. Sono guadagni di rappresentazione, non di
previsione, ed è opportuno tenerli distinti nella lettura dei risultati.

Va infine dichiarato, per la corretta interpretazione dei capitoli successivi, che le campagne
di misura ivi riportate sono state eseguite senza alcun raffinamento: griglia di ricerca
coincidente con quella dell'ambiente e sotto-celle coincidenti con le celle di ricerca. I
meccanismi descritti in questo capitolo appartengono dunque al sistema realizzato ma non sono
esercitati da quei dati, e nessuna delle differenze misurate è loro attribuibile. La
quantificazione del loro effetto sulle prestazioni resta fra gli sviluppi futuri.
